\section{Geographic CVD: Algorithm}
\label{sec:algorithm}

\subsection{Notation}

Let $G_v = (V_v, E_v)$ be the subgraph at bisection tree node $v$, with
$m = |V_v|$ tracts and population $p_i$ for each tract $i$.
Let $\ell_i = (\lambda_i, \phi_i)$ be the WGS84 (longitude, latitude) of the
population-weighted centroid of tract $i$'s polygon, loaded from the centroid
companion file (Section~\ref{sec:data-pipeline}).
Let $\mathbf{x}_i = (x_i, y_i) = \mathrm{proj}(\ell_i)$ be the Albers-projected
coordinates in meters (Section~\ref{subsec:albers}).

\subsection{Albers Equal-Area Conic Projection}
\label{subsec:albers}

We use the EPSG:5070 spherical Albers formula from Snyder~\citeyear{snyder1987},
Chapter~14:

\begin{definition}[Albers Projection --- Continental US (EPSG:5070)]
\label{def:albers}
Parameters: central meridian $\lambda_0 = -96^{\circ}$,
standard parallels $\phi_1 = 29.5^{\circ}$, $\phi_2 = 45.5^{\circ}$.
Given WGS84 $(\lambda, \phi)$ in degrees, convert to radians and compute:
\begin{align}
  n     &= \tfrac{1}{2}(\sin\phi_1 + \sin\phi_2), \\
  C     &= \cos^2\!\phi_1 + 2n\sin\phi_1, \\
  \rho_0 &= \sqrt{C} / n, \\
  \rho  &= \sqrt{C - 2n\sin\phi} / n, \\
  \theta &= n(\lambda - \lambda_0), \\
  x     &= R \cdot \rho\sin\theta, \\
  y     &= R(\rho_0 - \rho\cos\theta),
\end{align}
where $R = 6{,}371{,}000$\,m is the spherical Earth radius.
\end{definition}

\begin{remark}[Spherical approximation error]
Definition~\ref{def:albers} uses a spherical Earth ($R = 6{,}371{,}000$\,m).
The EPSG:5070 standard uses the GRS80 ellipsoid; the difference in
eccentricity ($e^2 = 0.00669438$) introduces at most 0.3\% geometric
error for continental US coordinates ($\lambda \in [-125^{\circ}, -66^{\circ}]$,
$\phi \in [24^{\circ}, 50^{\circ}]$).
At census-tract resolution (typical tract width 200--2{,}000\,m), this corresponds
to $\le 3$\,m positional error per tract.
For redistricting purposes, this error is negligible: it is far below the
minimum meaningful districting unit (a census tract) and introduces no
directional bias.
An ellipsoidal implementation using PROJ library bindings is deferred to
Phase~3.
\end{remark}

\begin{remark}[State-specific projections]
Alaska (FIPS 02) uses EPSG:3338 (Alaska Albers, central meridian $-154^{\circ}$,
parallels $55^{\circ}$/$65^{\circ}$); Hawaii (FIPS 15) uses EPSG:102007
(Hawaii Albers); Puerto Rico (FIPS 72) uses EPSG:32161.
Selection is automatic from the state FIPS code at load time.
The EPSG code used is recorded in the audit chain (\S\ref{subsec:audit}).
\end{remark}

\subsection{Seed Derivation}

Phase~2 uses a distinct seed prefix to isolate Phase~2 seeds from Phase~1
in the audit chain.

\begin{definition}[Geographic CVD Node Seed]
\label{def:geo-seed}
Let $\mathtt{base\_seed}$ be the unsigned 64-bit integer global seed and
$\mathtt{node\_path}$ the path string for bisection tree node $v$.
The \emph{geographic CVD node seed} is
\[
  s_v^{\mathrm{geo}} = \mathrm{first8}\bigl(
    \mathrm{SHA\text{-}256}(
      \texttt{"CVD\_GEO\_INIT\_"} \;\|\; \mathtt{node\_path}
      \;\|\; \texttt{"\_"} \;\|\; \mathrm{le8}(\mathtt{base\_seed})
    )
  \bigr),
\]
where $\mathrm{le8}(x)$ encodes $x$ as 8 little-endian bytes and
$\mathrm{first8}(h)$ reads the first 8 bytes as a little-endian
unsigned 64-bit integer~\citep{sha256nist}.
\end{definition}

The prefix \texttt{"CVD\_GEO\_INIT\_"} is strictly distinct from the Phase~1
prefix \texttt{"CVD\_INIT\_"}: a test invariant asserts both constants are
present in source and are not equal.

\subsection{K-Farthest Seed Initialisation (Euclidean)}

\begin{definition}[K-Farthest Euclidean Initialisation]
\label{def:geo-kfarthest}
Let $\mathrm{sorted} = (t_0,\ldots,t_{m-1})$ be tracts in $V_v$ ordered by
global index.
\begin{enumerate}
\item \textbf{First seed}: $\mathrm{seed}_0 = t_{s_v^{\mathrm{geo}} \bmod m}$.
\item \textbf{Euclidean distances from} $\mathbf{x}_{\mathrm{seed}_0}$:
  for each $t \in V_v$, $\delta(t) = \|\mathbf{x}_t - \mathbf{x}_{\mathrm{seed}_0}\|$.
\item \textbf{Second seed}: $\mathrm{seed}_1 = \arg\max_{t \in V_v} \delta(t)$
  (Euclidean farthest; ties broken by sorted index).
\end{enumerate}
The projected seed positions are $\mathbf{s}_0 = \mathbf{x}_{\mathrm{seed}_0}$,
$\mathbf{s}_1 = \mathbf{x}_{\mathrm{seed}_1}$.
\end{definition}

Unlike Phase~1 (BFS farthest), Definition~\ref{def:geo-kfarthest} finds the
two tracts that are geographically farthest apart in projected coordinates.
This is an $O(m)$ scan rather than a BFS pass.
For rectangular states, k-farthest Euclidean initialisation places seeds at
the two geographic corners, immediately producing near-converged Voronoi
regions.

\subsection{Voronoi Assignment and Centroid Update}

\begin{definition}[Euclidean Voronoi Assignment]
\label{def:euc-voronoi}
Given projected seeds $(\mathbf{s}_0, \mathbf{s}_1)$, the assignment
$a: V_v \to \{0,1\}$ assigns each tract to its nearest seed by Euclidean distance:
\[
  a(t) = \arg\min_{j \in \{0,1\}} \|\mathbf{x}_t - \mathbf{s}_j\|,
\]
with ties broken in favour of district 0.
\end{definition}

\begin{definition}[Population-Weighted Centroid Update]
\label{def:pw-centroid}
Let $D_j = \{t \in V_v : a(t) = j\}$ be the tracts in district $j$ and
$P_j = \sum_{t \in D_j} p_t$ the district population.
The \emph{population-weighted centroid} of $D_j$ in projected space is
\[
  \bar{\mathbf{c}}_j = \frac{1}{P_j}\sum_{t \in D_j} p_t \cdot \mathbf{x}_t.
\]
The point $\bar{\mathbf{c}}_j$ may not correspond to any real tract.
The updated seed position is the projected coordinate of the nearest real tract:
\[
  \mathbf{s}_j \leftarrow \mathbf{x}_{t^*},
  \quad t^* = \arg\min_{t \in V_v} \|\mathbf{x}_t - \bar{\mathbf{c}}_j\|.
\]
\end{definition}

The nearest-tract snap (Definition~\ref{def:pw-centroid}) is the only
approximation in Phase~2.
In continuous CVT theory, the seed would be exactly at $\bar{\mathbf{c}}_j$;
the snap replaces it with the geographically closest real tract.
This approximation is exact when $\bar{\mathbf{c}}_j$ lies inside a real tract's
footprint and introduces at most one census-tract width of error otherwise.

\subsection{$\varepsilon$-Convergence Criterion}

\begin{definition}[$\varepsilon$-Convergence]
\label{def:eps-convergence}
At the end of each iteration, let $\mathbf{s}_j^{\mathrm{new}}$ be the updated
seed and $\mathbf{s}_j^{\mathrm{prev}}$ the seed from the previous iteration.
The algorithm terminates when:
\[
  \max_{j \in \{0,1\}} \|\mathbf{s}_j^{\mathrm{new}} - \mathbf{s}_j^{\mathrm{prev}}\| < \varepsilon,
  \quad \varepsilon = 1.0\,\mathrm{m}.
\]
\end{definition}

The 1\,m tolerance is chosen to resolve discrete snap oscillations where two
equidistant candidate tracts alternate as the nearest-tract to
$\bar{\mathbf{c}}_j$.
Typical census tract widths are $\ge 200$\,m; a convergence tolerance of
1\,m is thus sub-tract-granularity and introduces no distributional bias.

\subsection{Formal Pseudocode}

Algorithm~\ref{alg:cvdgeo} gives the complete Phase~2 \cvdgeo{} bisection.

\begin{algorithm}[t]
\caption{\cvdgeo{}-\textsc{Bisect}($G_v$, $\mathbf{x}$, $\mathbf{p}$, $k_L$, $k_R$,
  $\mathtt{base\_seed}$, $\mathtt{node\_path}$, $\mathtt{cvd\_iters}$,
  $\mathtt{balance\_tolerance}$)}
\label{alg:cvdgeo}
\begin{algorithmic}[1]
\State $m \gets |V_v|$
\If{$m < 2$}
  \textbf{return} degenerate assignment
\EndIf
\State $s_v \gets \mathrm{CVD\text{-}Geo\text{-}Seed}(\mathtt{base\_seed}, \mathtt{node\_path})$
       \Comment{Definition~\ref{def:geo-seed}}
\State $\mathrm{sorted} \gets \mathrm{sort}(V_v)$;
       $\mathrm{seed}_0 \gets \mathrm{sorted}[s_v \bmod m]$
\State $\mathbf{s}_0 \gets \mathbf{x}_{\mathrm{seed}_0}$;
       $\mathbf{s}_1 \gets \mathbf{x}_{\arg\max_t \|\mathbf{x}_t - \mathbf{s}_0\|}$
       \Comment{k-farthest: Definition~\ref{def:geo-kfarthest}}
\State $\mathbf{s}_0^{\mathrm{prev}} \gets (-\infty, -\infty)$;
       $\mathbf{s}_1^{\mathrm{prev}} \gets (-\infty, -\infty)$
\For{$\mathrm{iter} = 0$ \textbf{to} $\mathtt{cvd\_iters} - 1$}
  \State $a(t) \gets \arg\min_{j} \|\mathbf{x}_t - \mathbf{s}_j\|$ for all $t \in V_v$
         \Comment{Voronoi: Definition~\ref{def:euc-voronoi}}
  \State $a \gets \mathrm{Rebalance}(a, G_v, \mathbf{p}, \mathtt{balance\_tolerance})$
         \Comment{200-iter boundary-swap}
  \For{$j \in \{0, 1\}$}
    \State $D_j \gets \{t : a(t) = j\}$;
           $P_j \gets \sum_{t \in D_j} p_t$
    \State $\bar{\mathbf{c}}_j \gets P_j^{-1}\sum_{t \in D_j} p_t \cdot \mathbf{x}_t$
           \Comment{pop-weighted centroid: Definition~\ref{def:pw-centroid}}
    \State $\mathbf{s}_j \gets \mathbf{x}_{\arg\min_{t \in V_v} \|\mathbf{x}_t - \bar{\mathbf{c}}_j\|}$
           \Comment{nearest-tract snap}
  \EndFor
  \If{$\max_j \|\mathbf{s}_j - \mathbf{s}_j^{\mathrm{prev}}\| < 1.0$}
    \textbf{break}
    \Comment{$\varepsilon$-convergence: Definition~\ref{def:eps-convergence}}
  \EndIf
  \State $\mathbf{s}_0^{\mathrm{prev}} \gets \mathbf{s}_0$;
         $\mathbf{s}_1^{\mathrm{prev}} \gets \mathbf{s}_1$
\EndFor
\State $a \gets \arg\min_j \|\mathbf{x}_t - \mathbf{s}_j\|$ for all $t$
       \Comment{final assignment from converged seeds}
\State $(S, \bar{S}) \gets \mathrm{Rebalance}(a, G_v, \mathbf{p}, \mathtt{balance\_tolerance})$
\State \textbf{return} $(S, \bar{S})$
\end{algorithmic}
\end{algorithm}

\paragraph{Key difference from Phase~1.}
The medoid computation (Phase~1: 100 BFS passes per iteration for the two
probe sets) is replaced by a weighted mean plus a nearest-neighbour scan
(Phase~2: $O(m)$ arithmetic per iteration with no BFS overhead).
All other structural elements --- k-farthest initialisation,
post-hoc rebalance, audit chain output --- are preserved.

\subsection{Convergence Analysis}

\begin{proposition}[Monotone Energy Descent for Phase~2]
\label{prop:geo-convergence}
Let $\mathcal{E}(\mathbf{s}_0, \mathbf{s}_1)
= \sum_{j} \sum_{t \in D_j} \|\mathbf{x}_t - \mathbf{s}_j\|^2$
be the total Euclidean quantisation energy.
Each iteration (Voronoi assignment + population-weighted centroid update +
nearest-tract snap) satisfies
$\mathcal{E}^{(i+1)} \le \mathcal{E}^{(i)} + O(\varepsilon^2)$.
\end{proposition}

\begin{proof}
(a) Voronoi assignment minimises $\mathcal{E}$ over all assignments given
    fixed seeds: assigning each tract to its nearest seed reduces
    $\sum_t \|\mathbf{x}_t - \mathbf{s}_{a(t)}\|^2$.
(b) Population-weighted centroid update minimises $\mathcal{E}$ over
    seed positions given fixed assignment: the minimiser of
    $\sum_{t \in D_j}\|\mathbf{x}_t - \mathbf{c}\|^2$ over $\mathbf{c}$ is the
    population-weighted mean $\bar{\mathbf{c}}_j$.
(c) Nearest-tract snap replaces the optimal $\bar{\mathbf{c}}_j$ with the closest
    real-tract projection $\mathbf{s}_j^* = \arg\min_{t}\|\mathbf{x}_t - \bar{\mathbf{c}}_j\|$.
    The additional energy from the snap is
    $\sum_{t \in D_j}\bigl(\|\mathbf{x}_t - \mathbf{s}_j^*\|^2 - \|\mathbf{x}_t - \bar{\mathbf{c}}_j\|^2\bigr)
    = 2(\mathbf{s}_j^* - \bar{\mathbf{c}}_j)^\top \sum_{t}(\mathbf{x}_t - \bar{\mathbf{c}}_j)
    + P_j\|\mathbf{s}_j^* - \bar{\mathbf{c}}_j\|^2$.
    The first term vanishes because $\bar{\mathbf{c}}_j$ is the centroid
    ($\sum_t p_t(\mathbf{x}_t - \bar{\mathbf{c}}_j) = 0$); the second term
    is $P_j\|\mathbf{s}_j^* - \bar{\mathbf{c}}_j\|^2$.
    The snap distance from the continuous centroid $\bar{c}$ to the nearest
    real tract $c_* = \arg\min_{t \in D_j} \|p_t - \bar{c}\|$ is bounded by
    the maximum nearest-neighbour radius
    $r_{\max} = \max_t \min_{t' \neq t} \|\mathbf{x}_t - \mathbf{x}_{t'}\|$
    of the projected tract set.
    For census-tract graphs with typical inter-centroid spacing $\geq 500\,$m,
    and convergence threshold $\varepsilon = 1\,$m, the snap displacement
    satisfies $\|c_* - \bar{c}\| \leq r_{\max}$, which is bounded by the
    minimum tract size.
    Thus the total energy change per snap is $O(r_{\max}^2)$, bounded and
    converging.
Combining (a), (b), (c): energy is non-increasing up to $O(\varepsilon^2)$
rounding per iteration.
Since $\mathcal{E} \ge 0$ and the number of distinct Voronoi partitions of
$V_v$ is finite, the iteration terminates.
\end{proof}

\subsection{Runtime Complexity}

\begin{proposition}[Phase~2 Runtime]
\label{prop:geo-runtime}
\cvdgeo{} bisection at a node with $m$ tracts requires:
\begin{itemize}
\item K-farthest initialisation: one Euclidean distance scan, $O(m)$.
\item Per iteration:
  \begin{itemize}
  \item Voronoi assignment: $O(m)$ (two Euclidean comparisons per tract).
  \item Centroid update: $O(m)$ weighted accumulation + $O(m)$ nearest scan.
  \end{itemize}
\item Post-hoc rebalance: $O(200 m)$.
\item Total: $O((n_{\mathrm{iter}} + 200) \times m)$.
\end{itemize}
For $n_{\mathrm{iter}} = 20$: $O(220m) = O(m)$.
This is strictly faster than Phase~1 at the same iteration count, which
requires $O(100m)$ BFS per iteration for the medoid computation in addition to
the $O(m)$ Voronoi assignment.
\end{proposition}

\subsection{Audit Chain Fields}
\label{subsec:audit}

Phase~2 appends the following fields to
\texttt{runs/\{label\}/\{year\}/index.json}
(new fields relative to Phase~1 are marked with $\dagger$):

\begin{verbatim}
"structure": "centroidal-voronoi",
"cvd_metric": "geographic",           (was "graph-distance")
"n_iter": 20,
"seed_init": "kmeans++",
"base_seed": 12345678,
"cvd_seed": 987654321,
"cvd_seed_formula": "SHA-256('CVD_GEO_INIT_'||node_path||'_'||base_seed:u64le)", (dagger)
"node_path": "01",                    (dagger)
"iterations_to_convergence": 11,
"convergence_warning": false,
"projection": "EPSG:5070",            (dagger)
"final_ec": 2847,
"final_pp_mean": 0.241
\end{verbatim}

The $\dagger$ fields are new in Phase~2.
\texttt{cvd\_seed\_formula} is a literal string encoding the seed derivation
so that an independent auditor can recompute \texttt{cvd\_seed} from
\texttt{base\_seed} and \texttt{node\_path} without consulting source code.
\texttt{node\_path} is required: the Phase~2 seed depends on both \texttt{base\_seed}
and \texttt{node\_path}, and without \texttt{node\_path} the seed cannot be
rederived from \texttt{base\_seed} alone.
\texttt{projection} records the EPSG code used; it varies by state FIPS code
(5070 for continental US; 3338 for Alaska; 102007 for Hawaii; 32161 for Puerto Rico).
